\part*{ВВЕДЕНИЕ}
\addcontentsline{toc}{part}{ВВЕДЕНИЕ}
В условиях быстрого развития технологий в последние несколько лет поддельные мультимедиа стали центральной проблемой. Используя передовые инструменты, создание поддельных медиапродуктов, таких как изображения и видео, становится очень простым, если доступны большие объёмы данных. Создание поддельного медиаконтента неэтично и представляет угрозу для общества. Они широко используются в киберпреступности для кражи личных данных, кибервымогательства, фейковых новостей, финансового мошенничества и многого другого. Поэтому больше внимания следует уделять выявлению контрафактной продукции. Обнаружение поддельного контента является серьёзной проблемой и пользуется большим спросом в сфере цифровой безопасности.

Особую сложность среди всех видов поддельных медиа представляют фальсифицированные изображения, а именно те случаи, когда изменён не весь кадр, а лишь отдельный фрагмент. Такие локальные манипуляции сложнее обнаружить, так как большая часть изображения остаётся подлинной. Современные генеративные модели способны создавать высокореалистичные фрагменты, визуально неотличимые от оригинала, что требует разработки новых, более совершенных алгоритмов обнаружения фальсификации.

Цель работы --- разработка и программная реализация метода автоматизированного определения области фальсифицированного фрагмента изображения с помощью двухпоточной свёрточной нейронной сети.

Для достижения поставленной цели необходимо выполнить следующие задачи:
\begin{itemize}[label = ---]
    \item провести обзор и сравнение существующих методов обнаружения фальсифицированного фрагмента на изображении;
    \item разработать метод обнаружения фальсифицированного фрагмента на изображении с использованием сверточной нейронной сети, привести IDEF0-диаграммы метода и схемы его алгоритмов;
    \item разработать программное обеспечение, реализующее данный метод;
    \item выбрать метрики для оценки качества работы метода и провести исследования качества работы метода.
\end{itemize}